\section{Regularity of equality sets}
\label{sec:equality-regularity}

The comparison established in \Cref{thm:main} turns an equality
set into a local perimeter almost-minimizer. We first identify its
regular boundary and prove the differentiability needed for the
curvature formulas. We then remove the singular set from the integral
identities that will locate the equality region.

\subsection{The perimeter support and its curvature}

\begin{proposition}[Boundary of an equality set]
\label{prop:equality-boundary}
Let $M^n$ be complete, simply connected and smooth, with $n\ge2$ and
$\Sec_M\le0$. Suppose that a bounded measurable set $E\subset M$ has
positive volume $V$, finite ambient perimeter $A$, and
\[
 A=n\omega_n^{1/n}V^{(n-1)/n}.
\]
Put $m=n-1$, $R=(V/\omega_n)^{1/n}$, and $H_0=m/R$. The support
$\Gamma=\supp|D\chi_E|$ is compact and has a decomposition
$\Gamma=\Sigma\mathbin{\dot\cup}Z$ with the following properties:
\begin{enumerate}
 \item $\Sigma$ is a relatively open, boundaryless, embedded
 $C^{1,1}_{\mathrm{loc}}$ hypersurface, and $Z$ is compact, with
 $\dim_{\HH}Z\le n-8$. In particular, $Z=\varnothing$ when $n<8$.
 \item The side occupied by $E$ defines an outward unit normal $N$
 on $\Sigma$, and
 \begin{equation}
 \label{eq:equality-perimeter-area}
 |D\chi_E|=d\sigma\lfloor\Sigma,
 \qquad \sigma(\Sigma)=A=s_mR^m,
 \qquad \overline\Sigma=\Gamma.
 \end{equation}
 Here $d\sigma$ denotes induced area on $\Sigma$, extended by zero
 to $M$.
 \item There are $C<\infty$ and $r_0>0$ such that
 \begin{equation}
 \label{eq:equality-area-growth}
 \sigma\bigl(\Sigma\cap B_r(x)\bigr)\le Cr^m
 \qquad(x\in\Gamma,\ 0<r<r_0).
 \end{equation}
 \item The outward trace mean curvature satisfies
 $H=\divg_{T\Sigma}N=H_0$ almost everywhere on $\Sigma$.
\end{enumerate}
These assertions concern the perimeter support and the measurable
class of $E$, not the topological boundary of an arbitrary
representative.
\end{proposition}

\begin{proof}
\emph{Almost-minimality and the regular part.}
Write $\theta=(n-1)/n$ and $\Lambda=A/V$. For every bounded
finite-perimeter set $F$, \Cref{thm:main} gives
\begin{equation}
\label{eq:equality-almost-minimality}
 \begin{split}
 P(E)-P(F)
 &\le A\left[1-\left(\frac{|F|}{V}\right)^\theta\right]
 \le\Lambda|E\triangle F|.
 \end{split}
\end{equation}
Indeed $z^\theta\ge\min\{z,1\}$ for $z\ge0$, and the difference
of the volumes is bounded by the volume of the symmetric difference.
If $E\triangle F$ is contained, up to a null set, in a compact subset
of a bounded open set $O$, locality of perimeter cancels the part
outside $O$. Thus
\[
 P(E;O)\le P(F;O)+\Lambda|E\triangle F|.
\]
It suffices to consider competitors with $P(F;O)<\infty$; agreement
with $E$ near the complement of $O$ then gives finite global
perimeter as well.

We apply the local smooth-manifold conclusion of
\cite[Corollary~1.6]{AntonelliPasqualettoPozzetta2022}. To check its
quasi-perimeter hypothesis, on measurable subsets of $O$ define
\[
 G(F)=\Lambda|F\triangle(E\cap O)|.
\]
This functional is finite at the empty set and satisfies
\[
 |G(F_1)-G(F_2)|\le\Lambda|F_1\triangle F_2|.
\]
Its continuity exponent is therefore $1>1-1/n$, as required in
Condition~(1.3) of that reference. The localized inequality says that
$E$ minimizes $P+G$ under compact modifications; in particular it
does so under the volume-preserving modifications of Definition~1.2.
For a ball $O$ centered on $\Gamma$, the remaining hypothesis
$P(E;O)>0$ follows from the definition of support. The cited local
conclusion applies in every dimension $n\ge2$ and requires no global
lower Ricci bound.

It gives a $C^{1,\alpha}_{\mathrm{loc}}$ regular part of full
perimeter, a relatively closed singular part of Hausdorff dimension
at most $n-8$, and local upper area growth. Here and below the local
exponent may be decreased so that $0<\alpha<1$. These are statements
about the perimeter support: one may take the density-one
representative supplied by the theorem, whose boundary equals that
support. For completeness, the perimeter measure is concentrated on
the reduced boundary, giving one inclusion. Conversely, a boundary
point of this representative is an essential boundary point; if the
perimeter vanished in a neighborhood, the characteristic function
would be constant almost everywhere there, a contradiction.

The support is unchanged by modifying $E$ on a null set. It is a
closed subset of a compact set containing $E$, since completeness
makes closed bounded sets compact. Hence $\Gamma$ and its closed
singular subset $Z$ are compact. The local area estimates become
\Cref{eq:equality-area-growth} with uniform constants by a finite
cover of $\Gamma$.

Near a regular point, the whole support is a graph dividing a
coordinate neighborhood into two connected sides. On either side
$D\chi_E=0$, so $\chi_E$ is constant almost everywhere. One way to
see this last fact is to express the distributional derivatives in
coordinates, compensating for the volume density, and mollify on
smaller coordinate balls. Overlapping balls give the same constant.
The two side values must differ: otherwise the characteristic
function would be constant across the neighborhood, since the graph
has zero ambient volume. The occupied side consequently defines a
consistent outward normal. The divergence formula on a graph
identifies its perimeter with its induced area. Together with the
full-perimeter conclusion, this proves
\Cref{eq:equality-perimeter-area}. In particular, every neighborhood
of a support point meets the regular part, so $\Sigma$ is dense in
$\Gamma$.

\emph{The graph equation.}
Choose graph coordinates $(y,z)\in\R^m\times\R$ with
$z=u(y)$ on $\Sigma$ and $E$ below the graph, up to a null set.
Let $G(y,z)$ be the ambient metric matrix and
$W(y,z)=\sqrt{\det G(y,z)}$. The graph area integrand is
\[
 L_g(y,z,p)
 =W(y,z)\left[(-p,1)^T G(y,z)^{-1}(-p,1)\right]^{1/2}.
\]
For a smooth compactly supported function $\zeta$ in the graph
domain, the graph $u+t\zeta$ gives admissible sets for both signs of
small $t$. The sharp deficit
$P(F)-n\omega_n^{1/n}|F|^\theta$ is nonnegative and vanishes at
$F=E$. Its volume derivative has coefficient
$n\omega_n^{1/n}\theta V^{\theta-1}=H_0$. Define
\[
 \mathcal B(y,z,p)=\partial_zL_g(y,z,p)-H_0W(y,z).
\]
Differentiating area and the volume between the two graphs gives
\begin{equation}
\label{eq:equality-graph-equation}
 \int\left[
  \partial_pL_g(y,u,Du)\cdot D\zeta
  +\mathcal B(y,u,Du)\zeta\right]dy=0.
\end{equation}
The derivative $D$ here is Euclidean differentiation in $y$.
The matrix $\partial_{pp}L_g$ is positive definite. Indeed, the
Hessian of a positive definite norm is positive in directions
transverse to its radial kernel, and a nonzero slope variation
$(-v,0)$ is not radial at $(-p,1)$. On compact sets of arguments the
equation is therefore uniformly elliptic.

The local graph-regularity result in \Cref{lem:equality-graph-regularity}
applies to this weak equation: a $C^{1,\alpha}$ graph stationary for
area minus $H_0$ times volume is locally $C^{1,1}$. Thus $Du$ is
locally Lipschitz. The proof of that result follows in the next
subsection; it obtains this conclusion without an advance bound on
second derivatives.

\emph{The signed curvature.}
We can now interpret the graph equation geometrically. On a
$C^{1,1}$ graph the induced metric and its volume density are locally Lipschitz.
For a tangential field $X=X^iF_i$ in a parametrization $F$, with
induced metric $(a_{ij})$, metric compatibility gives, almost
everywhere,
\[
 \divg_\Sigma X
 =\partial_iX^i+\tfrac12X^i a^{jk}\partial_i a_{jk}
 =\frac{1}{\sqrt{\det a}}\,
   \partial_i\bigl(\sqrt{\det a}\,X^i\bigr).
\]
The identities are also distributional: Lipschitz coefficients obey
the product rule, and the weak mixed derivatives of $F$ commute.
Consequently compactly supported tangential Lipschitz fields have
zero integral divergence. Splitting a field into its tangential
and normal parts now gives
\begin{equation}
\label{eq:equality-tangential-divergence}
 \divg_{T\Sigma}Y
 =\divg_\Sigma(Y^{\mathsf T})+H\langle Y,N\rangle,
 \qquad H=\divg_{T\Sigma}N,
\end{equation}
almost everywhere and weakly on each regular chart.

For the graph variation used in
\Cref{eq:equality-graph-equation}, take $Y=\zeta\,\partial_z$.
The area derivative is $\int\divg_{T\Sigma}Y\,d\sigma$, and the
volume derivative is its outward flux. By
\Cref{eq:equality-tangential-divergence}, the graph equation becomes
\[
 \int_\Sigma (H-H_0)\langle\partial_z,N\rangle
                       \zeta\,d\sigma=0.
\]
The factor $\langle\partial_z,N\rangle$ is strictly positive,
because $E$ occupies the lower side; equivalently its product with
the graph area density is $W$. Arbitrary compactly supported
$\zeta$ therefore give $H=H_0$ almost everywhere. The coefficient
$H_0$ came from the same global deficit on every patch, so this is
the same signed outward curvature on every component. This
completes the proof.
\end{proof}

\subsection{Local regularity of the graph equation}

The local input used above concerns the graph equation itself. Its
proof first bounds the energy of difference quotients, uniformly in
their step, and then controls their gradients by comparison with a
constant-coefficient equation.

\begin{lemma}[Regularity of a stationary graph]
\label{lem:equality-graph-regularity}
Let $\Omega\subset\R^m$ be open, with $m\ge1$, and let
$u\in C^{1,\alpha}_{\mathrm{loc}}(\Omega)$ for $0<\alpha<1$.
Let $G(y,z)$ be the matrix of a smooth Riemannian metric on a
neighborhood of the graph of $u$ in $\R^m\times\R$, and put
\[
 W(y,z)=\sqrt{\det G(y,z)},\qquad
 L_g(y,z,p)=W(y,z)
 \left[(-p,1)^TG(y,z)^{-1}(-p,1)\right]^{1/2}.
\]
Fix $H_0>0$ and write
$\mathcal B(y,z,p)=\partial_zL_g(y,z,p)-H_0W(y,z)$.
If
\[
 \int_\Omega\left[
 \partial_pL_g(y,u,Du)\cdot D\zeta
 +\mathcal B(y,u,Du)\zeta\right]dy=0
 \qquad(\zeta\in C_c^\infty(\Omega)),
\]
then $u\in C^{1,1}_{\mathrm{loc}}(\Omega)$.
\end{lemma}

\begin{proof}
\emph{Difference quotients and an energy bound.}
The positivity of $\partial_{pp}L_g$ was verified after
\Cref{eq:equality-graph-equation}. Shrink the graph chart so that
its arguments and the segments between nearby arguments lie in one
compact ellipticity region. Choose a Euclidean
ball $B_{4\rho}$ compactly contained in its parameter
domain. For a coordinate unit vector $e$ and $0<\delta<\rho$, put
\[
 v(y)=\frac{u(y+\delta e)-u(y)}{\delta},
 \qquad y\in B_{3\rho}.
\]
Each $v$ is $C^1$, though no uniform bound on $Dv$ is yet known.
The identity
\[
 v(y)=\int_0^1\partial_eu(y+s\delta e)\,ds
\]
does give a uniform $C^{0,\alpha}$ bound for $v$.

Subtract the translated equations and divide by $\delta$. To
describe the coefficients explicitly, set
\[
 \begin{aligned}
 T_s(y)&=(1-s)(y,u(y),Du(y))\\
       &\quad+s(y+\delta e,u(y+\delta e),Du(y+\delta e)).
 \end{aligned}
\]
The resulting weak equation is
\begin{equation}
\label{eq:equality-quotient-equation}
 \int\left[(aDv+f)\cdot D\zeta
              +(b\cdot Dv+d)\zeta\right]dy=0,
\end{equation}
where
\begin{align*}
 a&=\int_0^1\partial_{pp}L_g(T_s)\,ds,
 \qquad b=\int_0^1\partial_p\mathcal B(T_s)\,ds,\\
 f&=\int_0^1
       \bigl(\partial_e\partial_pL_g(T_s)
             +\partial_z\partial_pL_g(T_s)v\bigr)\,ds,\\
 d&=\int_0^1
       \bigl(\partial_e\mathcal B(T_s)
             +\partial_z\mathcal B(T_s)v\bigr)\,ds.
\end{align*}
In these expressions $\partial_e$ differentiates only the explicit
$y$ argument. The known $C^{1,\alpha}$ norm of $u$ bounds the
$C^{0,\alpha}$ norms of all $T_s$, uniformly in $s$ and $\delta$.
Smooth composition on the compact argument set, together with the
bound for $v$, shows that $a,f$ have uniform $C^{0,\alpha}$ bounds
and $b,d$ are uniformly bounded. The matrices $a$ are symmetric and
uniformly elliptic.

Test \Cref{eq:equality-quotient-equation} with $\chi^2v$, where
$\chi$ is supported in $B_{3\rho}$ and equals one on $B_{2\rho}$.
Such tests are justified by smooth approximation. Ellipticity gives
a multiple of $\int\chi^2|Dv|^2$ on the left. All other terms are
bounded by Young's inequality and the uniform bounds on
$v,f,b,d$, absorbing the terms linear in $Dv$ into that left side.
In particular, the first-order term involving $b\cdot Dv$ is
absorbed in this way. The remaining bound is
$C\int(\chi^2+|D\chi|^2)$, so
\begin{equation}
\label{eq:equality-quotient-energy}
 \int_{B_{2\rho}}|Dv|^2\,dy\le C,
\end{equation}
uniformly in the difference step. No bound on a second derivative of
$u$ has been used.

\emph{Harmonic comparison and the gradient bound.}
Fix $x\in B_\rho$ and $0<r\le\min\{\rho,1\}$. On $B_r(x)$,
let $w-v\in H^1_0(B_r(x))$ solve
\[
 \divg(a(x)Dw)=0.
\]
Existence and uniqueness follow by minimizing the positive quadratic
Dirichlet energy in $v+H^1_0(B_r(x))$. Put $z=v-w$. Subtract the
equations and test with $z$; the constant vector $f(x)$ integrates
to zero against $Dz$. With all norms on $B_r(x)$, ellipticity and
the coefficient bounds give
\[
 \begin{split}
 \lambda\|Dz\|_2^2
 &\le Cr^\alpha
      \bigl(\|Dv\|_2+|B_r|^{1/2}\bigr)\|Dz\|_2\\
 &\quad+C\bigl(\|Dv\|_2+|B_r|^{1/2}\bigr)\|z\|_2.
 \end{split}
\]
The second line includes the first-order term $b\cdot Dv$.
Zero-boundary Poincar\'e gives $\|z\|_2\le Cr\|Dz\|_2$.
Since $r\le r^\alpha$ for $r\le1$, it follows that
\begin{equation}
\label{eq:equality-frozen-comparison}
 \left(\fint_{B_r(x)}|Dv-Dw|^2\right)^{1/2}
 \le Cr^\alpha\left(
      1+\left(\fint_{B_r(x)}|Dv|^2\right)^{1/2}\right).
\end{equation}
Energy minimization after subtracting an affine function also gives,
for every constant vector $p$,
\begin{equation}
\label{eq:equality-harmonic-energy}
 \left(\fint_{B_r(x)}|Dw-p|^2\right)^{1/2}
 \le C\left(\fint_{B_r(x)}|Dv-p|^2\right)^{1/2}.
\end{equation}

We use the following elementary interior estimate for this constant
coefficient equation. For $0<\gamma\le1/4$,
\begin{equation}
\label{eq:equality-harmonic-oscillation}
 \left(\fint_{B_{\gamma r}(x)}
       |Dw-(Dw)_{B_{\gamma r}(x)}|^2\right)^{1/2}
 \le C\gamma
       \left(\fint_{B_r(x)}|Dw-p|^2\right)^{1/2}.
\end{equation}
Here subscripts on a function denote its average. To justify the
estimate at the weak regularity in use, mollify $Dw-p$ in an
interior ball. Its components solve the same constant coefficient
equation. A linear coordinate change, with distortion controlled by
ellipticity, makes them harmonic. The mean value formula reproduces
each such function by convolution with a smooth radial averaging
kernel of radius comparable to $r$. Differentiating that kernel
and applying Cauchy--Schwarz bounds its gradient on $B_{r/2}(x)$ by
\[
 Cr^{-1}\left(\fint_{B_r(x)}|Dw-p|^2\right)^{1/2}.
\]
The radius of the averaging kernel is chosen small enough that its
support remains in the transformed interior ball. Letting the
mollification scale tend to zero proves the same estimate for a
representative of $Dw-p$, and its oscillation on $B_{\gamma r}(x)$
is at most this bound times $2\gamma r$. This proves
\Cref{eq:equality-harmonic-oscillation}.

We control the mean gradient and its oscillation together: harmonic
comparison contracts the oscillation, while the coefficient errors
are summable on geometric scales. Define
\[
 p_r=\fint_{B_r(x)}Dv,
 \qquad
 J(r)=\left(\fint_{B_r(x)}|Dv-p_r|^2\right)^{1/2}.
\]
Using $p=p_r$ in the preceding estimates and comparing averages on
the two balls gives
\begin{align}
\label{eq:equality-oscillation-iteration}
 J(\gamma r)
 &\le C\gamma J(r)
   +C\gamma^{-m/2}r^\alpha\bigl(1+|p_r|+J(r)\bigr),\\
\label{eq:equality-mean-iteration}
 |p_{\gamma r}-p_r|&\le\gamma^{-m/2}J(r).
\end{align}
Fix $\gamma$ so that $C\gamma\le1/2$, and choose a fixed
$\ell\ge1$ with $\ell/2\ge\gamma^{-m/2}$. Then
\[
 S(r)=1+|p_r|+\ell J(r)
 \quad\hbox{satisfies}\quad
 S(\gamma r)\le(1+C'r^\alpha)S(r).
\]
Choose one starting radius $r_*\le\min\{\rho,1\}$. Its starting
averages are bounded uniformly for every $x\in B_\rho$ by
\Cref{eq:equality-quotient-energy}. On the geometric scales
$r_j=\gamma^jr_*$, the product of $1+C'r_j^\alpha$ is finite,
because $\sum_jr_j^\alpha<\infty$. Thus $|p_{r_j}|$ is uniformly
bounded. Each fixed-step quotient is $C^1$, so these averages tend
to $Dv(x)$ as $j\to\infty$. We obtain a bound for $Dv(x)$ uniform
in $x\in B_\rho$, in $e$, and in the sufficiently small step
$\delta$.

Finally,
\[
 Dv(y)=\frac{Du(y+\delta e)-Du(y)}{\delta}
\]
holds classically. The uniform bound just proved bounds increments
of $Du$ along every coordinate direction. On a smaller coordinate
cube, join two points by coordinate segments and add the increment
bounds. The sum of their lengths is at most $\sqrt m$ times the
Euclidean distance between the points. Therefore $Du$ is locally
Lipschitz and $u\in C^{1,1}_{\mathrm{loc}}(\Omega)$.
\end{proof}

\subsection{Integral identities across the singular set}

The regular boundary now has constant mean curvature. To use it in
the rigidity argument, we must justify integral tests that need not
vanish near its singular set.

\begin{lemma}[Cutoffs for an equality boundary]
\label{lem:equality-cutoffs}
Under the hypotheses of \Cref{prop:equality-boundary}, there are
$\eta_j\in\Lip_c(\Sigma)$ such that
\[
 0\le\eta_j\le1,
 \qquad \eta_j(x)\longrightarrow1\quad(x\in\Sigma),
 \qquad
 \int_\Sigma|\nabla_\Sigma\eta_j|^2\,d\sigma\longrightarrow0.
\]
\end{lemma}

\begin{proof}
If $Z=\varnothing$, then $\Sigma=\Gamma$ is compact, and we take
$\eta_j=1$. Otherwise $n\ge8$, $m\ge7$, and
\[
 \dim_{\HH}Z\le m-7<m-2,
 \qquad \HH^{m-2}(Z)=0.
\]
The compactness, finite area, embedding and uniform growth
hypotheses of \Cref{lem:capacity-cutoffs} follow from
\Cref{prop:equality-boundary}. That lemma supplies the stated
cutoffs. In particular their supports are compact in the manifold
topology of $\Sigma$, and the construction does not assume that
$\Sigma$ is complete or has finitely many components.
\end{proof}

For every smooth ambient field $Y$ defined near $\Gamma$, the
cutoffs and \Cref{eq:equality-tangential-divergence} imply
\begin{equation}
\label{eq:global-first-variation}
 \int_\Sigma\divg_{T\Sigma}Y\,d\sigma
 =H_0\int_\Sigma\langle Y,N\rangle\,d\sigma.
\end{equation}
Indeed, apply the compact-support identity to $\eta_jY$ and expand:
\begin{align*}
 &\int_\Sigma\eta_j\divg_{T\Sigma}Y\,d\sigma
 +\int_\Sigma\langle\nabla_\Sigma\eta_j,Y^{\mathsf T}\rangle\,d\sigma\\
 &\hspace{35mm}=H_0\int_\Sigma\eta_j\langle Y,N\rangle\,d\sigma.
\end{align*}
The middle integral has absolute value at most
$\|Y\|_\infty A^{1/2}\|\nabla_\Sigma\eta_j\|_2$, which tends
to zero. All other terms converge by dominated convergence, since
the ambient data are bounded near compact $\Gamma$. The same
identity holds on any compact boundaryless component of $\Sigma$
directly, by integrating
\Cref{eq:equality-tangential-divergence} on that component.

When $n\ge4$, the same cutoffs extend
\Cref{thm:extrinsic-sobolev} to every smooth ambient function $v$
restricted to $\Sigma$. Here the signed normalized mean curvature is
$h=H/m=1/R$, and we use the case $b=0$. With $q=2m/(m-2)$ and $c=4/(m-2)$, the result is
\begin{equation}
\label{eq:equality-sobolev}
 c\int_\Sigma|\nabla_\Sigma v|^2\,d\sigma
 +\frac{m}{R^2}\int_\Sigma v^2\,d\sigma
 \ge m s_m^{2/m}
       \left(\int_\Sigma|v|^q\,d\sigma\right)^{2/q}.
\end{equation}
To justify passage from $\eta_jv$ to $v$, boundedness of $v$ and
its gradient near $\Gamma$ gives convergence in $L^2$ and $L^q$,
and
\[
 \nabla_\Sigma(\eta_jv)-\nabla_\Sigma v
 = (\eta_j-1)\nabla_\Sigma v
    +v\nabla_\Sigma\eta_j
 \longrightarrow0\quad\hbox{in }L^2(\Sigma).
\]
The first term tends to zero by dominated convergence and the
second by \Cref{lem:equality-cutoffs}. Thus the whole Dirichlet
energy, including its cross term, converges. In particular the
constant function and fixed ambient perturbations of it are now
legitimate tests; no completeness of the regular locus has been
introduced.
